\documentclass[12pt]{article}
\usepackage[utf8]{inputenc}
\usepackage[T1]{fontenc}
\usepackage{hyperref}
\title{Yuki Language Resources for K-12 Curriculum Development}
\author{Compiled by Claude Sonnet 5.5 (claude-sonnet-5-5) for Andr\'as Kornai}
\date{}
\begin{document}
\maketitle

\section{Introduction}
Yuki (ISO 639-3 \texttt{yuk}, Glottocode yuki1243) is a Yukian language of the Round Valley area of northern California, related to Wappo (Yuki-Wappo family). Search results give its dialects as Round Valley Yuki, Huchnom and Coast Yuki. They report the last native speaker, Arthur Anderson, as dying in 1983. The language is therefore dormant, and any curriculum would be heritage and revitalization work. One Wikipedia page says revival work is ongoing, with the language taught in grade school alongside Wailaki on the Round Valley Reservation. I could not confirm this from a primary source, so treat it with caution.

\section{Primary Educational Resources}
I found no primers, textbooks or classroom materials for Yuki.

\section{Linguistic and Reference Materials}
\begin{itemize}
\item Wikipedia, ``Yuki language''. It mentions Kroeber's 1911 sketch and Sawyer and Schlichter's 1984 lexical corpus. It also mentions Balodis's 2016 \emph{Yuki Grammar}. \url{https://wikimili.com/en/Yuki_language}
\item Wikipedia, ``Northern Yukian language''. \url{https://en.wikipedia.org/wiki/Northern_Yukian_language}
\item Wikipedia, ``Yuki people''. \url{https://en.wikipedia.org/wiki/Yuki_people}
\item East Carolina University library catalogue preview, returned by the search for Yuki materials. I did not confirm which item it describes. \url{https://lib.ecu.edu/catalog-preview/catalog/2397601}
\end{itemize}

\section{Religious and Community Resources}
I found no scripture or community organisation resources.

\section{Digital and Audio Resources}
The American Philosophical Society holds Yuki materials in the Harvey Pitkin Papers, mostly Kroeber's stories and texts. \url{https://indigenousguide.amphilsoc.org/entry/38366}
I found no audio or video resources.

\section{Curriculum Development Context}
There are no living first-language speakers. A curriculum would have to be built from archival sources, namely Kroeber's texts, the Sawyer and Schlichter lexical corpus and Balodis's grammar. Community partners on the Round Valley Reservation would need to lead the work, and I could not verify any existing programme.

\end{document}